\clearpage\thispagestyle{empty}
\pdfbookmark[0]{Quellen zu Komet (15. Februar 1823)}{cometfeb1823-sources}
\begin{center}\textbf{Quellen zu Komet (15. Februar 1823)}\end{center}
\noindent Diese KI-Übertragung ins heutige Deutsch folgt dem Bericht \textit{Comet} in \href{https://gdz.sub.uni-goettingen.de/id/PPN236007467}{Gauß’ \textit{Werken}, Bd.~VI} (Göttingen, 1874), S.~442--443. Der Erstdruck erschien unter der Überschrift \textit{Göttingen} in den \textit{Göttingischen gelehrten Anzeigen}, Nr.~26, 15. Februar 1823, S.~\href{https://images.p.sub.uni-goettingen.de/iiif/image/gdz:PPN385030444_1823_1:00000251/full/1600,/0/default.jpg}{249}, \href{https://images.p.sub.uni-goettingen.de/iiif/image/gdz:PPN385030444_1823_1:00000252/full/1600,/0/default.jpg}{250} und \href{https://images.p.sub.uni-goettingen.de/iiif/image/gdz:PPN385030444_1823_1:00000253/full/1600,/0/default.jpg}{251}. Die Digitalisate stammen aus der Niedersächsischen Staats- und Universitätsbibliothek Göttingen. Der Bericht gibt Beobachtungen wieder, die Rümker Gauß in einem Brief aus Paramatta mitteilte. Er verweist auf den \hyperlink{native-425}{Bericht vom 24. Mai 1819} über Enckes Entdeckung.
\par\smallskip\noindent Der Erstdruck schreibt Enckes Namen \textit{Enke} und verwendet andere Zeitzeichen und Dezimalzeichen. Dort fehlen der eingeklammerte Verweis auf S.~420 und die Punktreihen neben den beiden Bedeckungszeiten. Die Tabelle und die Angaben zur Bedeckung und zur Länge folgen der Werkausgabe, einschließlich der nur stellenweise angegebenen Einheiten, der Schreibung \textit{Paramatta} und der Sternzeitstunde 13 am 22. Juni, die auch in der Zeitschrift steht. Der Bericht lässt offen, was \textit{mittlere} bei den Koordinaten bedeutet. KI-Übertragung und KI-Satz.
\clearpage
